% ============================================================================
% Ch1 ground-up rebuild (W2, 2026-08-02): drafted from the frozen sources per
% DISSERTATION_ARCHITECTURE.md (Ch1) and AUDIT_DOSSIER.md. Aim/primer/universe
% prose adapted from the June ch1 only where dossier Part 3 marks it safe;
% every number carries a same-line evidence comment per the blueprint's
% verification protocol.
% W3 assembly (2026-08-02): prelude.tex's "items submitted" table now carries the
% five deliverables of Section 1.3, item for item — verified during assembly.
% ============================================================================
\chapter{Introduction}
\label{ch:intro}

\section{Project Aim}
\label{sec:intro-aim}

The aim of this project is to build the two instruments that tell apart two explanations of an apparent forecasting gain.
When a text model appears to beat a price-based benchmark at forecasting volatility, the filing may have carried genuine
information about future risk, or the model may have recognised the filer and recalled how volatile that firm usually is.

The first is a survivorship-free, point-in-time benchmark linking sixteen years of United States regulatory disclosures to the stock-market volatility their filers subsequently realised. The second is an identity-controlled audit that separates what a filing \emph{says} from \emph{which firm filed it}. A benchmark alone cannot tell the two apart, and an audit without a leakage-free benchmark has nothing trustworthy to measure. The project therefore delivers both instruments, returns a calibrated verdict on the incremental forecasting value of disclosure text, and releases them so that the same audit can be rerun wherever entities repeatedly write text about themselves.

Both instruments rest on two financial ingredients, stated plainly here for a computer science reader and developed formally in Chapter~\ref{ch:background}. The \emph{volatility} of a stock is how widely its price swings from day to day. \emph{Realised volatility} is the measurable, after-the-fact version of that quantity: here, the annualised standard deviation of a firm's daily returns over the 5, 10 or 20 trading days following a filing. % src: frozen 03_dataset.tex / 04_methods.tex (horizons, target)
Volatility forecasts size trading positions, price option contracts and set regulatory capital \cite{andersen2003modeling,poon2003forecasting}. Unlike the direction of returns, volatility is genuinely forecastable, because calm and turbulent periods cluster in time \cite{bollerslev1986garch,engle1982arch}. A three-coefficient regression of future volatility on the price history's own recent behaviour, the HAR-RV model \cite{corsi2009har}, is consequently notoriously hard to beat. Any claim that text helps must be measured against that standard rather than against the weak baselines much of the applied literature prefers. The price history is not the only competitor. Markets impound public information almost immediately, so a text signal must survive the benchmark and a market that has already read the filing.

The candidate text is the regulatory disclosure record. Companies listed in the United States must report to the Securities and Exchange Commission (SEC), whose EDGAR system publishes every filing with a precise timestamp \cite{sec2024edgar}. A substantial literature reports risk-relevant signal in filing word features, finance-specific dictionaries and media tone \cite{kogan2009predicting,loughran2011liability,tetlock2007giving}, and modern text encoders have renewed the expectation \cite{beltagy2020longformer,devlin2019bert}. Two disclosure channels are studied because they differ in length, cadence and content. \emph{Long-form} periodic reports (the annual 10-K and quarterly 10-Q) are book-length accounts of a firm's business and risks, tens of thousands of tokens long and beyond the context window of any standard encoder. \emph{Event-driven} 8-K current reports are short and must be filed within days of a material event. \textbf{SP500Vol-Bench}, the benchmark built here, aligns 144{,}129 such filings (31{,}601 long-form and 112{,}528 event-driven) to the volatility outcomes of their filers under an audited no-look-ahead rule, % src: frozen 03_dataset.tex; TECHNICAL_SUPPLEMENT.md Table S1
yielding 431{,}245 aligned filing-by-horizon rows over 2010--2025. % src: frozen 03_dataset.tex (rows derive from the 144,129 aligned filings)
Its point-in-time S\&P~500 universe holds 914 membership intervals covering 832 securities and retains every firm later delisted or acquired. % src: crsp_restricted/membership_intervals.csv (verified 2026-08-01)

% Aim-level justification of the material (assessor point: why this dataset).
% Conceptual case only; operational reasons (liquidity, EDGAR coverage, the CRSP
% licence) are argued in Chapter 3. No numbers, hence no src comments needed.
Two features recommend this material over more convenient text. The first is that the record is mandated rather than volunteered. Periodic reports fall due on a fixed calendar whether or not the period was eventful, and a current report is triggered by rule, not by an editor's judgement of interest. The firm writes the account itself, not an intermediary, and the moment it becomes public is a matter of record, which is what makes the benchmark's point-in-time claim auditable. The second is that volatility is a target a disclosure could plausibly inform, as the direction of returns is not. Forecasting direction from a public document requires the market to have mispriced it; forecasting volatility asks only how wide the range of outcomes will be, a quantity the text must improve upon relative to what the price history already implies, and one an enumerated risk narrative is written to describe. The mandate cuts both ways. Compelled prose hedges and repeats, so any signal is thin and diffuse, as Chapter~\ref{ch:background} details. The regime rests on the premise that a firm's own account of its risks is useful to those bearing them, and a calibrated ceiling on what it adds beyond price measures one corner of that premise. A null answer is worth as much here as a positive one.

The audit half of the aim answers a threat that the field's usual precautions leave untouched. A company that was turbulent last year is usually turbulent this year, and every company here files again and again under its own name and in its own house style. A model expressive enough to notice that it has read this filer many times before can therefore score well on filings it has never seen without learning anything from them, recalling instead how risky that firm normally is. Splitting the data by date does not stop it, because the training years' filers are largely the test years' filers too. Outside finance this has a name: shortcut learning \cite{geirhos2020shortcut}. On financial text it is documented already. A baseline given a firm's identity and none of its words matches models that read its earnings calls \cite{yu2025samecompany}. This project therefore asks a harder question than the customary one: not whether text beats a price baseline, but whether it beats a reference already told which firm is filing. Where it does not, the project asks how much of the apparent advantage was the name rather than the news.

The verdict, reached in Chapters~\ref{ch:results} and~\ref{ch:validation}, is stated here because the whole report is organised around it: a calibrated near-null with one bounded exception. The near-null is an ordering, and the report keeps its three claims apart: disclosure text on this panel is detectable, only partially attributable to content, and not realisable as economic value. Two units carry the counts. A \emph{comparison} scores one challenger against the price baseline at one horizon on one channel, asking whether text stands alone; a \emph{combination cell} pairs a text model with a price reference on the same terms, asking whether it adds. Standalone, no challenger beats the price baseline in any of 180 squared-error comparisons, and 155 are significantly worse. % src: results/tables/dm_pairwise_clustered.md
 % src: results/tables/variance_unit_standalone180.csv (0 of 153 under squared error, volatility-unit and variance-unit QLIKE, raw and Holm); TECHNICAL_SUPPLEMENT.md S7
Against a single \emph{recalibrated} price reference, 38 of 69 combination cells show an apparent, placebo-confirmed increment; Figure~\ref{fig:intro-ladder} follows it through the stronger references that dissolve it, beside what each rung could have detected had a signal of known size been there. % src: results/tables/m1_ensemble_primary.md; results/tables/control_intersection_ensemble.md
 % src: results/tables/firm_identity_ensemble.md; results/tables/maximal_reference_ensemble.md; results/tables/control_intersection_ensemble.md
 % src: results/tables/signal_injection_power.csv
An anonymisation arm re-scores all 112{,}528 event-driven filings with their identifying language masked. It prices the identity share of the apparent 8-K gains at 0.51 (pooled median; 95 per cent interval $[0.44, 0.56]$). % src: results/tables/anon_arm.md; results/tables/anon_share_ci.md
A bounded residual survives. Prompted large-language-model readings of 8-K filings retain 0.2--0.5 per cent over the firm-identity reference, % src: results/tables/firm_identity_ensemble.md (+0.52/+0.24/+0.21%)
and deliver no significant portfolio value in any of that arm's six economic cells. % src: results/tables/portfolio_econ.csv (0 of C6's 6; grid-wide 1 of 18)
The audit also travels. On an external earnings-call corpus it reprices a published text gain to nothing detectable (0 of 24 cells), % src: results/tables/maec_audit.md
while on consumer reviews the identity control is itself costly and a content residual survives. % src: results/tables/yelp_multiarm.md
The apparatus detects signal where signal exists, and the shortcut's size is a property of the panel, not a constant.

\begin{figure}[tbp]
\centering
\includegraphics[width=\textwidth]{figures/C1_reference_ladder.pdf}
\caption[The reference ladder and its power calibration]{The reference ladder, read left to right as the reference is told more. Each rung re-runs the same 69 combination cells and no rung reads any text. R2 and R3 strengthen R1 in different directions, not in sequence; their survivor sets are disjoint cell by cell, not merely by channel, which is why R4, requiring both, is empty. R3's pool is HAR, SHAR, GARCH, EGARCH and ARIMA. The 153 at the top left are the census's text and fusion comparisons. Each count gives the number of cells whose text increment survives Holm correction~\cite{holm1979simple}, and at R1 the label-shuffle placebo as well; R4's count is zero because no cell survives.} % caption src: Chapters~\ref{ch:data} and~\ref{ch:results}; counts as cited in this section
\label{fig:intro-ladder}
\end{figure}

That verdict rests on a wider panel than the proposal named. The firm universe, framed at scoping around the Nasdaq-100, was refined with the supervisor's agreement to the point-in-time S\&P~500 membership held in the CRSP records \cite{crsp2026data}, whose audited survivorship-free provenance and roughly five-fold larger cross-section the question demands. % src: frozen 03_dataset.tex (universe change); dossier Part 1.1
Chapter~\ref{ch:data} documents the change and its engineering in full. The project was scoped so that a rigorous, reproducible answer constitutes success whatever its sign, and the chapters that follow argue that the carefully bounded near-null above is exactly such an answer.

\section{Objectives}
\label{sec:intro-objectives}

Building both instruments and returning that verdict decomposes the aim into five objectives. The first three restate, at the final supervisor-agreed scope, the dataset, modelling and evaluation commitments of the Scoping and Planning document. The fourth and fifth extend them, added as the identity threat moved to the centre of the investigation. Chapter~\ref{ch:conclusions} returns a verdict on each.

\begin{enumerate}
  \item \textbf{Benchmark construction (O1).} Construct SP500Vol-Bench: a survivorship-free, point-in-time benchmark linking the regulatory disclosures of S\&P~500 constituents over 2010--2025 to forward realised volatility at the 5-, 10- and 20-day horizons. Construction uses point-in-time entity linking, chronological splits, and an automated no-look-ahead audit over the full aligned corpus.
  \item \textbf{Model spectrum (O2).} Implement and evaluate a model spectrum wide enough that the verdict cannot be blamed on a weak representative: price-only baselines, classical text pipelines, fine-tuned Transformer encoders under five long-document strategies, frozen seven-to-eight-billion-parameter embedding models, and a prompted 32-billion-parameter language model. % src: frozen 04_methods.tex (model matrix)
  Fine-tuned models share one fixed recipe and three replicate seeds, one tuned arm apart, so that differences measure representation rather than tuning effort or run-to-run luck. % src: frozen 05_protocol.tex (seeds 2026--2028)
  \item \textbf{Incremental-value audit protocol with power calibration (O3).} Design and execute an incremental-value audit that credits text only for the out-of-sample loss it removes beyond a \emph{recalibrated} price reference, under significance tests that treat the trading day as the independent unit, multiplicity control within pre-declared families, label-shuffle placebos, and a power calibration that reports the minimum detectable effect beside every null verdict.
  \item \textbf{Identity attribution (O4).} Separate content from identity in whatever survives. Audit every apparent gain against identity-aware references, then price its identity share by direct intervention: anonymising filings, swapping documents between matched firms, and probing models with zero-content prompts.
  \item \textbf{Portability (O5).} Establish that the audit, and not only its verdict on this panel, is the contribution. Run the protocol end-to-end on two external public panels (consumer-review ratings and earnings-call volatility), and report what it detects, what it reprices, and what it leaves standing.
\end{enumerate}

Each objective has a site of delivery: O1, the data layer of Chapter~\ref{ch:data} (Section~\ref{sec:data-sources}); O2, the model matrix of Section~\ref{sec:design-models} and the standalone census of Chapter~\ref{ch:results}; O3, the design of Section~\ref{sec:design} and its execution through Chapters~\ref{ch:results} and~\ref{ch:validation}; O4, the ladder and identity instruments of Chapters~\ref{ch:results} and~\ref{ch:validation}; O5, the external-panel evidence of Chapter~\ref{ch:validation}.

\section{Deliverables}
\label{sec:intro-deliverables}

The Scoping and Planning document committed the project to a reproducible dataset, implemented models, an evaluated incremental-value analysis and this report. Those commitments are delivered in the five artefacts below; the audit protocol and the run-fingerprint chain have grown into deliverables in their own right. Appendix~\ref{app:external} describes how each is accessed; the first four are shared through the public repository at \url{https://github.com/Beiciccc/sp500vol-bench}.

\begin{enumerate}
  \item \textbf{SP500Vol-Bench, with a public core.} The benchmark of Chapter~\ref{ch:data}, released as the accession index of all 144{,}129 benchmark filings with timestamps, effective trading days and split assignments; % src: frozen 12_reproducibility.tex; release/DATA_CARD.md
  the 914 survivorship-free membership intervals and the 30{,}100 point-in-time link windows are withheld with the market data, being CRSP-derived, and are rebuilt by the ingestion script; % src: crsp_restricted/README.md; scripts/ingest_wrds.py
  the chronological split definition and a datasheet; and a licence-free variant that rebuilds prices from public sources, its coverage priced honestly in Chapter~\ref{ch:data}. Market-data-derived per-row values are withheld under licence (Section~\ref{sec:intro-ethics}), and their regeneration is fully scripted.
  \item \textbf{The audit protocol and its power calibration.} The reference ladder, identity controls and inference stack of Chapter~\ref{ch:data}, together with the registered analysis records that pre-declared its test families, the aggregate evidence tables behind every claim, and the injected-signal calibration that states what the apparatus could have detected. Other benchmark builders can port the protocol, as Chapter~\ref{ch:validation} demonstrates on two external panels.
  \item \textbf{The run-fingerprint chain.} Cryptographic configuration fingerprints of all 240 production runs (233 distinct digests), with a released script that re-hashes and verifies them, % src: results/tables/config_fingerprints.md
  tying every archived result to the exact code and settings that produced it.
  \item \textbf{The code repository.} The full regeneration pipeline (ingestion, point-in-time alignment, labelling, training, evaluation, figures and tables) under version control, guarded by 186 automated tests % src: writing/dissertation/supporting/REPRODUCIBILITY.md
  and by generators that abort if a source table drifts from its recorded values.
  \item \textbf{This report.}
\end{enumerate}

\section{Legal, Social, Ethical and Professional Issues}
\label{sec:intro-ethics}

The project involves no human participants, but its legal, social, ethical and professional dimensions were each assessed systematically and several shaped the design directly; Appendix~\ref{app:ethics} records the full assessment. The unit of analysis is the firm. The corpus is public corporate disclosures and licensed market data; the executive names filings carry are located only to be masked, by the anonymisation arm of Chapter~\ref{ch:results}, and no personal attribute is extracted or analysed. Ethically, no individual-level inference of any kind is drawn, so the University's ethical-review procedures for human participants and personal data are not engaged, a conclusion reached by documented assessment against the personal-data definitions of the UK General Data Protection Regulation and the Data Protection Act 2018, not by assumption. Legally, lawful access was honoured at both sources: EDGAR is public but rate-limited, and the download pipeline respects the SEC's fair-access policy, declares its identity and caches every response so that no document is requested twice \cite{sec2024edgar}. CRSP market data, accessed through Wharton Research Data Services, are licensed for research use and may not be redistributed \cite{crsp2026data}. That licence determines the shape of the release. Returns, labels, price features and the predictions built on them are withheld; the benchmark index and the aggregate evidence tables are released; the crosswalks carry CRSP identifiers and dates and are withheld with them. Any licence holder regenerates the withheld rows exactly from the released scripts. Socially, overstated predictive claims in finance can mislead investors and lend unearned authority to fragile systems. Every model in this report is a research artefact built to answer a scientific question, not a trading system and not investment advice. Professionally, in a literature visibly tilted towards positive findings, the project treats the honest reporting of a carefully bounded, largely null verdict as itself an obligation.

\section{Structure of the Report}
\label{sec:intro-structure}

The chapters run in the order the verdict was earned: standard, instruments, result, then the attempt to break it. Chapter~\ref{ch:background} sets the standard and justifies the methods adopted; Chapter~\ref{ch:data} builds both instruments and the engineering that makes the programme verifiable; Chapter~\ref{ch:results} runs them, from the standalone census to the economic pricing; Chapter~\ref{ch:validation} attacks the result from five directions; and Chapter~\ref{ch:conclusions} returns a verdict on each objective, states the project's limitations without deletion, and identifies the genuinely open future work. Two conventions hold throughout. Where the main text states a headline count, the appendix table or figure carrying it cell by cell is named beside it; and the arm, strategy and basis terms are defined in the Nomenclature.
